% Options for packages loaded elsewhere
\PassOptionsToPackage{unicode}{hyperref}
\PassOptionsToPackage{hyphens}{url}
%
\documentclass[
  11pt,
]{article}
\usepackage{amsmath,amssymb}
\usepackage{iftex}
\ifPDFTeX
  \usepackage[T1]{fontenc}
  \usepackage[utf8]{inputenc}
  \usepackage{textcomp} % provide euro and other symbols
\else % if luatex or xetex
  \usepackage{unicode-math} % this also loads fontspec
  \defaultfontfeatures{Scale=MatchLowercase}
  \defaultfontfeatures[\rmfamily]{Ligatures=TeX,Scale=1}
\fi
\usepackage{lmodern}
\ifPDFTeX\else
  % xetex/luatex font selection
\fi
% Use upquote if available, for straight quotes in verbatim environments
\IfFileExists{upquote.sty}{\usepackage{upquote}}{}
\IfFileExists{microtype.sty}{% use microtype if available
  \usepackage[]{microtype}
  \UseMicrotypeSet[protrusion]{basicmath} % disable protrusion for tt fonts
}{}
\makeatletter
\@ifundefined{KOMAClassName}{% if non-KOMA class
  \IfFileExists{parskip.sty}{%
    \usepackage{parskip}
  }{% else
    \setlength{\parindent}{0pt}
    \setlength{\parskip}{6pt plus 2pt minus 1pt}}
}{% if KOMA class
  \KOMAoptions{parskip=half}}
\makeatother
\usepackage{xcolor}
\usepackage[margin=1in]{geometry}
\usepackage{longtable,booktabs,array}
\usepackage{calc} % for calculating minipage widths
% Correct order of tables after \paragraph or \subparagraph
\usepackage{etoolbox}
\makeatletter
\patchcmd\longtable{\par}{\if@noskipsec\mbox{}\fi\par}{}{}
\makeatother
% Allow footnotes in longtable head/foot
\IfFileExists{footnotehyper.sty}{\usepackage{footnotehyper}}{\usepackage{footnote}}
\makesavenoteenv{longtable}
\setlength{\emergencystretch}{3em} % prevent overfull lines
\providecommand{\tightlist}{%
  \setlength{\itemsep}{0pt}\setlength{\parskip}{0pt}}
\setcounter{secnumdepth}{-\maxdimen} % remove section numbering
\ifLuaTeX
  \usepackage{selnolig}  % disable illegal ligatures
\fi
\usepackage{bookmark}
\IfFileExists{xurl.sty}{\usepackage{xurl}}{} % add URL line breaks if available
\urlstyle{same}
\hypersetup{
  hidelinks,
  pdfcreator={LaTeX via pandoc}}

\author{}
\date{}

\begin{document}

\section{Counterfactual RH search on the original split-supported theta
input}\label{counterfactual-rh-search-on-the-original-split-supported-theta-input}

\subsection{Scope}\label{scope}

The preceding Deligne-bound audit is not a no-go theorem for geometry
over F1, for the monoidal base tau, or for a future purity construction
using the full arithmetic input. Its stated conclusion concerns a
particular inference from an auxiliary exponential family's Frobenius
purity to a bound for the original theta metric. Its examples do not
assert that their chosen polynomials divide the original function
\(g=2\xi\).

This continuation takes the counterexample search as an independent
research objective. It proves a finite-sign certificate for a
counterexample to RH, directly from derivatives of the original theta
function. It also proves that an actual counterexample would eventually
supply such a certificate. The finite calculation reported below found
no negative witness. No counterexample to RH is claimed.

The power-sum/moment method belongs to established zero-positivity
methods; see Zhang (2015). The derivation below gives all its maps,
signs, and constants for the present source, rather than claiming a new
general RH criterion.

\subsection{1. Original input and
support}\label{original-input-and-support}

Retain

\[
G(R)=\{\tau\}\sqcup R^\bullet,\qquad e_R=0_R^\bullet,
\qquad p_R(\tau)=0,\quad p_R(r^\bullet)=r,
\]

and the coefficient square

\[
\begin{array}{ccc}
G(\mathbb Z)&\longrightarrow&G(\mathbb C)\\
p_{\mathbb Z}\downarrow&&\downarrow p_{\mathbb C}\\
\mathbb Z&\longrightarrow&\mathbb C.
\end{array}
\]

The support observation sends \(\tau\) to Boolean zero and every
represented coefficient to Boolean one. On the original support-labelled
coefficient fibres, a linear map lifts as \((\ell,v)\mapsto(\ell,fv)\),
with external absence fixed. A killed coefficient therefore reaches
\((\ell,0)\), not external absence. No conclusion about the sign of an
arithmetic functional follows just from that support value.

Use the unchanged source

\[
\Theta\phi(x)=2\sum_{n\ge1}\phi(nx),\qquad
\phi_*(x)=(4\pi^2x^4-6\pi x^2)e^{-\pi x^2},
\]

\[
f_0=\Theta\phi_*,\qquad
\mathcal Mf_0(s)=g(s)=2\xi(s)
=s(s-1)\pi^{-s/2}\Gamma(s/2)\zeta(s).
\]

These are the arithmetic-input equations (A1)--(A3), also identified
with the specified Hochschild trace in
\texttt{HOCHSCHILD\_COMPARISON.md}, equations (H4)--(H10). Define the
explicit logarithmic-coordinate function

\[
\psi(y)=e^{y/2}f_0(e^y).
\tag{1}
\]

Mellin substitution gives

\[
g(1/2+iz)=\int_{\mathbb R}\psi(y)e^{izy}\,dy.
\tag{2}
\]

The original reflection \(f_0(x)=x^{-1}f_0(1/x)\) gives
\(\psi(-y)=\psi(y)\). This reflection also follows by applying the
source Fourier identity to \(\phi_*\): Fourier sends \(D\) to \(1-D\),
so it fixes \((D^2-D)e^{-\pi x^2}=\phi_*\).

For \(x\ge1\), every summand in \(f_0(x)\) is positive, since

\[
4\pi^2n^4x^4-6\pi n^2x^2
=2\pi n^2x^2(2\pi n^2x^2-3)>0.
\]

Reflection proves positivity for all \(x>0\). The Gaussian series also
gives a constant \(C\) such that

\[
0<\psi(y)\le C\exp(9|y|/2-\pi e^{2|y|}).
\tag{3}
\]

For \(y\ge0\), bound the convergent sum of \(n^4e^{-\pi(n^2-1)e^{2y}}\)
by its value at \(y=0\); reflection supplies the other half-line. Thus
all differentiations and all moments below converge absolutely.

\subsection{2. Coefficients with the original mass
retained}\label{coefficients-with-the-original-mass-retained}

Put

\[
M_{2j}=\int_{\mathbb R}y^{2j}\psi(y)\,dy
=g^{(2j)}(1/2),\qquad
 a_j=\frac{(-1)^jM_{2j}}{(2j)!}.
\tag{4}
\]

Define the entire function \(F\) by its specified coefficients:

\[
F(w)=\sum_{j\ge0}a_jw^j,
\qquad F(z^2)=g(1/2+iz),\qquad a_0=g(1/2)>0.
\tag{5}
\]

No square-root branch is part of the definition, and \(a_0\) is not
assigned the value one. The map \(z\mapsto z^2\) identifies the two
reflected coordinates only at this scalar-function observation; the
original packet modules are retained in Section 5.

Define \(b_n\) by the analytic logarithmic derivative near zero:

\[
-\frac{F'(w)}{F(w)}=\sum_{n\ge0}b_nw^n.
\tag{6}
\]

The complete finite coefficient map is

\[
\boxed{
 b_n=\frac{-(n+1)a_{n+1}-\sum_{j=1}^{n}a_jb_{n-j}}{a_0}.
}
\tag{7}
\]

It follows by multiplying (6) by \(F\) and extracting the coefficient of
\(w^n\). For example,

\[
 b_0=-\frac{a_1}{a_0},\qquad
 b_1=\frac{a_1^2}{a_0^2}-\frac{2a_2}{a_0},\qquad
 b_2=-\frac{a_1^3}{a_0^3}+\frac{3a_1a_2}{a_0^2}-\frac{3a_3}{a_0}.
\tag{8}
\]

This operation need not preserve positivity of arbitrary coefficient
sequences. The matrices used below are not positive by their definition.

\subsection{3. A finite certificate contradicting
RH}\label{a-finite-certificate-contradicting-rh}

Let \(\alpha_j\) be the distinct zeros of \(F\), with their
multiplicities \(m_j\), and put \(\beta_j=\alpha_j^{-1}\). These symbols
are used to prove the test, not supplied to its numerical
implementation.

Estimate (3), followed by \(t=e^{2|y|}\), bounds the maximum modulus of
\(g\) by an exponential of \(O(R\log R)\). The maximum modulus of \(F\)
at radius \(R\) is consequently bounded by an exponential of
\(O(\sqrt R\log R)\). Its order is at most \(1/2<1\). Genus-zero
Hadamard factorization therefore gives

\[
F(w)=a_0\prod_j(1-\beta_jw)^{m_j},
\qquad \sum_jm_j|\beta_j|<\infty.
\tag{9}
\]

There is no nonconstant exponential factor of integer degree at this
order. Differentiating locally uniformly away from zeros proves

\[
\boxed{b_n=\sum_jm_j\beta_j^{n+1}.}
\tag{10}
\]

The exact coordinate of a zeta zero \(\rho=1/2+\delta+i\gamma\) is

\[
 z=-i(\rho-1/2)=\gamma-i\delta,\quad
 \alpha=z^2=\gamma^2-\delta^2-2i\gamma\delta,
\]

\[
\boxed{
 \beta=-\frac1{(\rho-1/2)^2}
 =\frac{\gamma^2-\delta^2+2i\gamma\delta}
 {(\gamma^2+\delta^2)^2}.
}
\tag{11}
\]

The two reflected roots \(\rho,1-\rho\) give one zero \(\alpha\) of the
same multiplicity in \(F\). Conjugation gives the conjugate \(\alpha\)
and \(\beta\).

For every real \(s\), (1)--(3) give
\(g(s)=\int\psi(y)e^{(s-1/2)y}dy>0\). Hence \(F\) has no nonpositive
real zeros: \(F(-v^2)=g(1/2-v)>0\). Thus RH is equivalent to every
\(\beta_j\) being positive real.

For \(d\ge0\), define the actual Hankel coefficient matrix

\[
\mathsf H_d=(b_{r+s})_{0\le r,s\le d}.
\tag{12}
\]

For \(P(X)=\sum_{r=0}^d c_rX^r\) with real coefficients, the exact value
is

\[
\boxed{
 c^T\mathsf H_dc
 =\sum_jm_j\beta_jP(\beta_j)^2.
}
\tag{13}
\]

Absolute convergence follows from (9) and boundedness of the polynomial
on the compact set of reciprocal zeros and zero. Conjugate pairs make
(13) real.

Under RH, each term is nonnegative. Therefore

\[
\boxed{
 c^T\mathsf H_dc<0\quad\Longrightarrow\quad\text{RH is false}.
}
\tag{14}
\]

The certificate may use \(c\in\mathbb Q^{d+1}\): a strict negative value
persists under sufficiently small rational coefficient perturbations.
With validated enclosures \(b_n\in[\underline b_n,\overline b_n]\),
exact rational interval arithmetic on \(\sum c_rc_sb_{r+s}\) is enough
to certify that its upper endpoint is negative.

A second necessary family is \(\mathsf H_d^{(1)}=(b_{r+s+1})\). Under RH
its quadratic form is \(\sum_jm_j\beta_j^2P(\beta_j)^2\ge0\). We tested
both, though the unshifted family alone is complete by the next
construction.

\subsection{4. Work through the off-line counterfactual: construct its
negative
direction}\label{work-through-the-off-line-counterfactual-construct-its-negative-direction}

Assume for this argument that RH is false. Equation (11) supplies a
nonreal reciprocal zero \(\beta\) and its conjugate, each of
multiplicity \(m\). They are nonreal because \(g\) has no real zeros, so
both \(\gamma\) and \(\delta\) are nonzero.

Put \(r=|\beta|/2\). There are finitely many distinct reciprocal zeros
of modulus at least \(r\). Remove \(\beta,\overline\beta\) from that
finite set and call the remaining set \(\mathcal A\). It is closed under
conjugation. Define the real polynomial

\[
B(X)=\prod_{\gamma\in\mathcal A}(X-\gamma).
\tag{15}
\]

Choose \(c\in\mathbb C\) satisfying \(\beta c^2=-1\); either of the two
choices works. For an integer \(N\ge0\), put

\[
w_N=\frac{c}{\beta^NB(\beta)},\qquad
v_N=\frac{\operatorname{Im}w_N}{\operatorname{Im}\beta},\qquad
u_N=\operatorname{Re}w_N-\operatorname{Re}\beta\,v_N,
\]

\[
\boxed{P_N(X)=X^NB(X)(\nu_N+v_NX).}
\tag{16}
\]

The coefficients are real. At the two target points,

\[
P_N(\beta)=c,\qquad P_N(\overline\beta)=\overline c.
\]

At every other point of modulus at least \(r\), \(P_N\) vanishes. The
two target contributions to (13) are exactly

\[
m\beta c^2+m\overline\beta\,\overline c^{\,2}=-2m.
\tag{17}
\]

Here is an explicit bound on everything else. Let

\[
B_0=\prod_{\gamma\in\mathcal A}(r+|\gamma|),\qquad
C=\frac{|c|B_0}{|B(\beta)|}
\left(1+\frac{|\operatorname{Re}\beta|+r}{|\operatorname{Im}\beta|}\right),
\]

\[
T=\sum_{|\beta_j|<r}m_j|\beta_j|<\infty.
\]

For \(|X|\le r\), the coefficient definitions imply

\[
|P_N(X)|\le C\left(\frac r{|\beta|}\right)^N=C2^{-N}.
\]

Consequently the complete, real quadratic form satisfies

\[
\boxed{
\mathcal L(P_N^2)
:=\sum_jm_j\beta_jP_N(\beta_j)^2
\le -2m+C^2T4^{-N}.
}
\tag{18}
\]

For all sufficiently large \(N\) the right side is strictly negative.
Every omitted reciprocal zero is included in the tail estimate, and
every multiplicity is retained. This proves

\[
\boxed{\mathrm{RH}\ \Longleftrightarrow\ \mathsf H_d\succeq0
\text{ for every }d\ge0.}
\tag{19}
\]

It also explains the limitations of a finite search. The degree needed
in (16) depends on the location of the off-line pair, the other larger
reciprocal zeros, and the tail. The proof does not say that a
counterexample must be visible in a small matrix. The polynomial (16) is
a counterfactual construction; no actual off-line pair has been used to
instantiate it.

\subsection{5. Connect the witness to the original finite packet,
including its
jets}\label{connect-the-witness-to-the-original-finite-packet-including-its-jets}

For a finite conjugation- and reflection-stable packet of actual zeros,
retain the original inclusion

\[
\sigma_h:E_h=\mathbb C[s]/(h)\hookrightarrow Q,
\qquad D_Q\sigma_h=\sigma_hA_h.
\]

The original arithmetic input constructs this through the kernel of
multiplication by \(j_hg\). An arbitrary polynomial example is not
automatically such a packet.

Because \(g(1/2)>0\), the operator \(A_h-\tfrac12I\) is invertible on
every actual packet. Define the specified rational functional calculus

\[
T_h=-(A_h-\tfrac12I)^{-2},\qquad
\mathbb C[X]\longrightarrow\operatorname{End}_{\mathbb C}(E_h),
\quad P\longmapsto P(T_h).
\tag{20}
\]

On a full local block \(A_h=\rho I+N_\rho\), \(N_\rho^m=0\), the matrix
is exactly

\[
\boxed{
T_h=\sum_{j=0}^{m-1}(-1)^{j+1}(j+1)
(\rho-1/2)^{-j-2}N_\rho^j.
}
\tag{21}
\]

The coefficient of \(N_\rho\) is \(2(\rho-1/2)^{-3}\ne0\), so this local
coordinate change retains the nilpotent length. The two reflected
primary modules can have the same \(T_h\) eigenvalue, but they remain
separate in \(E_h\).

The packet contribution to (13) is

\[
\boxed{\mathcal L_h(P^2)=\tfrac12
\operatorname{Tr}_{E_h}\bigl(T_hP(T_h)^2\bigr).}
\tag{22}
\]

The factor \(1/2\) counts each reflected pair once, matching (5). The
trace records algebraic multiplicities, while the full operators and
their nilpotent coefficients remain upstream of the trace map. Absolute
convergence in (9) makes the full functional the limit over packets
exhausting the actual zero set.

The existing critical-line comparison satisfies

\[
\ker(\Gamma\sigma_h)=E_{h,\mathrm{off}}.
\]

Equation (22) is evaluated on the source packet before applying that
map. An off-critical class maps to \((\ell,0)\) under the split lift of
\(\Gamma\), but is not deleted from the source on which the trace and
negative witness are evaluated. This is why the search does not obtain a
false positive or a false proof by restricting to the critical-only
target.

\subsection{6. The performed search uses the original theta seed, not
zero
ordinates}\label{the-performed-search-uses-the-original-theta-seed-not-zero-ordinates}

The numerical calculation uses (4), in the explicit form

\[
M_{2j}=4\sum_{n\ge1}\int_0^\infty
 y^{2j}\left(4\pi^2n^4e^{9y/2}-6\pi n^2e^{5y/2}\right)
 e^{-\pi n^2e^{2y}}\,dy.
\tag{23}
\]

The factor four consists of the source theta factor two and the
reflection factor two. The implementation computes
\(M_0,\ldots,M_{64}\), forms (7), and tests both matrix families through
dimension 16 by an unpivoted \(LDL^T\) calculation in their original
monomial coordinates. No rescaling or preconditioning is used. No zero
ordinate is supplied.

Two runs were completed:

\begin{longtable}[]{@{}
  >{\raggedright\arraybackslash}p{(\columnwidth - 8\tabcolsep) * \real{0.1765}}
  >{\raggedleft\arraybackslash}p{(\columnwidth - 8\tabcolsep) * \real{0.2353}}
  >{\raggedleft\arraybackslash}p{(\columnwidth - 8\tabcolsep) * \real{0.2353}}
  >{\raggedright\arraybackslash}p{(\columnwidth - 8\tabcolsep) * \real{0.1765}}
  >{\raggedright\arraybackslash}p{(\columnwidth - 8\tabcolsep) * \real{0.1765}}@{}}
\toprule\noalign{}
\begin{minipage}[b]{\linewidth}\raggedright
Working decimal precision
\end{minipage} & \begin{minipage}[b]{\linewidth}\raggedleft
Gauss--Legendre nodes
\end{minipage} & \begin{minipage}[b]{\linewidth}\raggedleft
Retained theta terms
\end{minipage} & \begin{minipage}[b]{\linewidth}\raggedright
y interval
\end{minipage} & \begin{minipage}[b]{\linewidth}\raggedright
Outcome
\end{minipage} \\
\midrule\noalign{}
\endhead
\bottomrule\noalign{}
\endlastfoot
120 & 384 & 12 & {[}0,4{]} & All 16 pivots positive numerically, in both
families \\
180 & 768 & 14 & {[}0,4{]} & Same outcome \\
\end{longtable}

Selected values are

\[
a_0=0.9942415563766282198255474793707954396\ldots,
\]

\[
b_0=0.0231049931154189707889338104303390140\ldots,
\qquad
b_1=0.00003717259928526968616486626247174057845\ldots,
\]

\[
b_2=0.0000001441739314009732796953815560948209\ldots,
\]

\[
\det\mathsf H_1
=b_0b_2-b_1^2
=1.949335554819142214739878219172139506\times10^{-9}\ldots.
\]

The largest relative difference between the two runs' stored pivot
values is below \(3\times10^{-79}\). Those stored values contain 80
significant digits, so the comparison does not support a stronger
precision claim. The scripts report the input choices and complete pivot
sequences.

These are exploratory multiple-precision quadratures, not
interval-certified integrations. A failed attempt to install an
interval-arithmetic backend did not produce an interval certificate;
only the documented mpmath computations are reported. Positive numerical
pivots do not prove the corresponding signs, and a negative numerical
pivot would require certified analytic enclosures before it could count
as a disproof. No negative numerical pivot was found.

\subsection{7. Finite verification and the next decisive
object}\label{finite-verification-and-the-next-decisive-object}

Six standard-library exact rational checks passed in normal and
optimized Python. They verify the logarithmic recurrence, positive-node
identity, an explicit nonreal-pair negative control, and a rational
interval evaluator that cannot call an interval containing zero strictly
negative. The negative control is expressly an auxiliary polynomial, not
an actual zeta packet.

The interval evaluator's contract is limited: it checks the quadratic
form from supplied valid intervals. Establishing that those intervals
enclose the actual \(b_n\) remains an analytic task. No Lean certificate
is claimed.

The decisive counterexample object in this branch is now a finite tuple

\[
(d,c_0,\ldots,c_d,
[\underline b_0,\overline b_0],\ldots,
[\underline b_{2d},\overline b_{2d}])
\]

with rational \(c_r\), rigorously justified coefficient enclosures from
the original theta input, and a strictly negative upper bound for
\(\sum_{r,s}c_rc_sb_{r+s}\). Equation (14) proves exactly what that
tuple would establish. Equation (18) proves that every actual off-line
counterexample has a finite witness of this type at some degree.

The calculation performed here did not find that tuple. The failure of
the previous auxiliary-purity inference is not used as evidence that one
exists. Geometry over the original absolute base remains a separate
possible source of further constraints; this note establishes no
impossibility theorem about it.

\subsection{Source and execution
record}\label{source-and-execution-record}

\begin{itemize}
\tightlist
\item
  Original arithmetic source: supplied \texttt{arithmetic\_input.tex},
  equations (A1)--(A10), including the actual multiplication-by-\(j_hg\)
  kernel sequence.
\item
  Original support reconstruction and internal homology: supplied
  SplitZero reconstruction record, Sections 1--4. The source-labelled
  zero and the external absent value are kept distinct throughout.
\item
  Original critical comparison: supplied
  \texttt{HOCHSCHILD\_COMPARISON.md}, equations (H16)--(H25). The
  visible text supplies the exact off-critical kernel and critical-jet
  retractions.
\item
  Prior scope: supplied \texttt{Tau\_Deligne\_Bound\_Audit/NOTE.tex},
  ``Result and scope''. It explicitly does not refute the split-zero
  construction or establish falsity of the arithmetic estimate for
  \(2\xi\).
\item
  Classical context: R. Zhang, \emph{On Power Sums of Positive Numbers},
  arXiv:1510.03420 (2015). Its abstract documents moment-theoretic
  positivity criteria for genus-zero entire functions and applications
  to RH. The concrete criterion and separation proof above are proved
  here, not attributed as a verbatim theorem from an uninspected full
  text.
\item
  Classical background: Hadamard factorization for entire functions of
  order less than one; the functional equation and completed-zeta
  notation are also recorded in NIST DLMF Section 25.4. The source
  Gaussian estimate above supplies the order bound needed for this
  application.
\item
  Computation: mpmath 1.3.0. Results are in \texttt{probe\_120dps.json}
  and \texttt{probe\_180dps.json}; reproduction script is
  \texttt{search\_theta\_moments.py}. Exact finite checks and the
  interval-input verifier are in \texttt{exact\_witness\_checks.py}.
\end{itemize}

No GitHub file, branch, comment, or remote publication was changed in
this continuation.

\end{document}
